\documentclass[11pt,a4paper]{article}
\usepackage[T1]{fontenc}
\usepackage[utf8]{inputenc}
\usepackage[main=italian,provide=*]{babel}
\usepackage{amsmath,amssymb}
\usepackage{geometry}
\geometry{margin=2.5cm}
\usepackage{booktabs}
\usepackage{seqsplit}
\usepackage{hyperref}
\hypersetup{colorlinks=true,linkcolor=blue,citecolor=blue,urlcolor=blue}

\title{VQE senza termine DM — trimero a catena aperta}
\author{Note di lavoro — Tesi triennale, Samuele Schiavi\\ Relatore: Prof.\ Alessandro Chiesa}
\date{}

\begin{document}
\maketitle

\section*{Introduzione}

Questo documento raccoglie i file della fase VQE \textbf{senza} termine DM per il
trimero a catena aperta — mirror diretto, per $N=3$, di quanto fatto per il dimero
(\texttt{vqe\_dimer.py}) e per l'anello (documento gemello). Un ansatz PMA con
branching e blocco RBS (produzione principale), lo stesso con blocco $W$, un ansatz
PMA senza branching, due notebook "spiegato" (guide interattive) e un notebook di
confronto sistematico.

\textbf{Scope.} Tutti i file qui descritti evitano deliberatamente il termine DM. Per
la catena, il VQE con DM è ancora completamente da fare — le scelte qui descritte non
dipendono da come verrà affrontato in futuro. Per i moduli \texttt{.py}, oltre a
descrizione/verifiche/conclusione, sono elencate anche le funzioni interne e come si
usa il file.

\section*{\texttt{vqe\_trimer\_chain.py} + \texttt{vqe\_trimer\_chain.png}}

Mirror diretto di \texttt{vqe\_trimer\_ring.py}, costruito da zero in questa sessione:
ansatz HA (9 parametri) e PMA con RBS e branching sui soli due bond fisici $(1,2)$ e
$(2,3)$. Per $J>0$ il fondamentale a basso campo è sempre il blocco $B$ (mai serve
scegliere fra $B$ e $C$ come nell'anello): $2K+3$ parametri, 5 a $K=1$ — contro i 6
dell'anello, un blocco RBS in meno.

\textbf{Verifiche.} Self-test a precisione macchina ($1.9\times10^{-15}$), script
principale completo eseguito: $\mathcal F=1$ esatto su tutti i 10 punti, per entrambi
gli ansatz. Figura ispezionata visivamente.

\textbf{Conclusione.} Modulo di produzione pronto, stesso standard di rigore
dell'anello.

\textbf{Funzioni interne.}
\begin{itemize}
\item \texttt{make\_ha\_ansatz(reps)} — l'ansatz hardware-efficient generico, stesso
ruolo della versione dell'anello.
\item \texttt{rbs\_block(qc,\ phi,\ q0,\ q1)} — il blocco RBS, identico a quello
dell'anello.
\item \texttt{\_select\_target\_block(J,\ b)} — sceglie \texttt{(block, M)} del
fondamentale: sempre $B$ sotto $b_c$ per $J>0$ (nessun confronto $E_B$ vs $E_C$
necessario, a differenza dell'anello), $A$ sopra.
\item \texttt{\_prepare\_initial\_state(qc,\ block,\ M)} — banale (solo $X$) per
$A$ estremo, generico (\texttt{prepare\_state}) per $B$ — qui $B$ non è mai banale
(a differenza del blocco a basso campo tipico dell'anello).
\item \texttt{make\_pma\_ansatz(J,\ b,\ K)} — l'ansatz PMA: stato iniziale + $K$
layer di RBS sui \emph{due} bond fisici + $R_y$ indipendenti finali, $2K+3$
parametri.
\item \texttt{\_energy(params,\ ansatz,\ hamiltonian,\ estimator)} — la funzione
obiettivo per \texttt{\seqsplit{scipy.optimize.minimize}}.
\item \texttt{run\_vqe(J,\ b,\ ansatz\_type,\ reps,\ K,\ n\_restarts,\ seed)} — VQE
completo a un punto (nessun \texttt{dm\_mode}/\texttt{D}: qui lo scope esclude il DM
fin dalla firma della funzione, non solo dal suo uso).
\item \texttt{vqe\_sweep(b\_values,\ J,\ ansatz\_type,\ reps,\ K,\ n\_restarts,\
seed,\ verbose)} — chiama \texttt{run\_vqe} su una griglia di $b$.
\item \texttt{\_self\_test()} — verifica l'ansatz a parametri nulli contro lo stato
di Kambe atteso.
\item \texttt{\_plot\_row}, \texttt{plot\_grid} — stessa funzione della versione
dell'anello, per la figura a 2$\times$4 pannelli.
\end{itemize}

\textbf{Come si usa.} Come libreria: \texttt{from vqe\_trimer\_chain import
run\_vqe, vqe\_sweep, \ldots}. Come script standalone: \texttt{python3
vqe\_trimer\_chain.py} esegue il self-test, poi lo sweep completo per HA e PMA e
salva \texttt{vqe\_trimer\_chain.png}/\texttt{.pdf}. Dipende da
\texttt{trimer\_chain\_exact.py} nella stessa cartella.

\section*{\texttt{vqe\_trimer\_chain\_W.py} + \texttt{vqe\_trimer\_chain\_W.png}}

Mirror di \texttt{vqe\_trimer\_ring\_W.py}: gate $W$, branching, $2K$ parametri (2 a
$K=1$, contro i 3 dell'anello).

\textbf{Verifiche.} Self-test a precisione macchina, script principale eseguito:
$\mathcal F=1$ esatto ovunque, con errore energetico esattamente nullo in più punti
sopra $b_c$.

\textbf{Conclusione.} Completa, insieme al file precedente, la coppia RBS/$W$ già
disponibile per l'anello.

\textbf{Funzioni interne.} Stessa struttura di \texttt{vqe\_trimer\_chain.py}, con
queste differenze:
\begin{itemize}
\item \texttt{w\_block(qc,\ theta,\ q0,\ q1)} — il blocco $W$ (CNOT-$R_y$-CNOT) al
posto di \texttt{rbs\_block}.
\item \texttt{make\_pma\_ansatz(J,\ b,\ K)} — solo $K$ blocchi $W$ sui due bond,
$2K$ parametri, nessun layer finale di $R_y$.
\item \texttt{make\_pma\_ansatz\_forced(branch,\ K)} — ramo iniziale fissato
indipendentemente da $b$; per il ramo \texttt{"basso"} usa
\texttt{prepare\_state} sullo stato di Kambe $(B,-\tfrac12)$ (diversamente
dall'anello, dove il ramo "basso" forzato usa la preparazione banale del
singoletto).
\item \texttt{run\_vqe\_forced(J,\ b,\ branch,\ K,\ n\_restarts,\ seed)} — il VQE
con il ramo forzato.
\item Le altre funzioni (\texttt{make\_ha\_ansatz}, \texttt{\_select\_target\_block},
\texttt{\_prepare\_initial\_state}, \texttt{\_self\_test}, \texttt{\_energy},
\texttt{run\_vqe}, \texttt{vqe\_sweep}, \texttt{\_plot\_row}, \texttt{plot\_grid})
hanno lo stesso ruolo della versione RBS.
\end{itemize}

\textbf{Come si usa.} Come libreria: \texttt{from vqe\_trimer\_chain\_W import
run\_vqe, vqe\_sweep, \ldots}. A differenza di \texttt{vqe\_trimer\_chain.py}, in
questa sessione è stato eseguito tramite uno script esterno equivalente (stesso
schema self-test + sweep + plot). Dipende da \texttt{trimer\_chain\_exact.py}.

\section*{\texttt{vqe\_trimer\_chain\_nobranch.py}}

Mirror di \texttt{vqe\_trimer\_ring\_nobranch.py}: stato iniziale fisso, RBS ripetuto
ad ogni ciclo sui due bond fisici, $5K$ parametri (contro i $6K$ dell'anello).

\textbf{Verifiche e scoperta.} Trovato e corretto un bug reale durante la
costruzione: il conteggio dei parametri per ciclo era sbagliato (4 invece di 5),
causava un \texttt{IndexError} immediato. Dopo la correzione, verificato un risultato
non presunto per analogia: a $b=1.0$ (sotto $b_c$), $K=1$ (5 parametri) resta vicino al
tetto strutturale $\approx5/6$ già noto per la famiglia "esatta" (RBS fatto una sola
volta), ma $K=2$ (10 parametri, RBS ripetuto due volte) lo \textbf{supera},
raggiungendo $\mathcal F=1$ esatto — coerente col pattern già visto per l'anello, dove
la famiglia ciclica a $K=2$ recupera l'esatto sotto DM dove $K=1$ resta bloccato.

\textbf{Conclusione.} Ripetere il blocco RBS (non solo aggiungere $R_y$) è ciò che
apre nuovo spazio raggiungibile — un'osservazione strutturale, non specifica del DM,
che vale anche nel regime $D=0$ qui documentato.

\textbf{Funzioni interne.}
\begin{itemize}
\item \texttt{rbs\_block(qc,\ phi,\ q0,\ q1)} — stesso blocco RBS già validato in
\texttt{vqe\_trimer\_chain.py}.
\item \texttt{make\_ha\_ansatz(reps)} — l'ansatz hardware-efficient, non usato dal
resto del modulo ma incluso per uniformità con gli altri file.
\item \texttt{pma\_2q\_trimer\_nobranch(K)} — l'ansatz senza branching: $X$ fisso
su un qubit, poi $K$ cicli di [RBS sui \emph{due} bond + $R_y$ indipendenti sui tre
qubit], $5K$ parametri (era scritto erroneamente come $4K$ prima della correzione).
\item \texttt{\_energy(params,\ ansatz,\ hamiltonian,\ estimator)} — stessa
funzione obiettivo degli altri moduli.
\item \texttt{run\_vqe(J,\ b,\ K,\ n\_restarts,\ seed)} — VQE a un punto, interfaccia
analoga agli altri moduli ma senza \texttt{ansatz\_type}/\texttt{reps}.
\item \texttt{vqe\_sweep(b\_values,\ J,\ K,\ n\_restarts,\ seed,\ verbose)} — chiama
\texttt{run\_vqe} su una griglia di $b$.
\end{itemize}

\textbf{Come si usa.} Come libreria: \texttt{from vqe\_trimer\_chain\_nobranch
import run\_vqe, vqe\_sweep}. Nessun blocco \texttt{\_\_main\_\_}, nessuna funzione
di self-test dedicata — la verifica più diretta è chiamare \texttt{run\_vqe} su un
punto qualsiasi e controllare \texttt{fidelity}. Dipende da
\texttt{trimer\_chain\_exact.py}.

\section*{\texttt{vqe\_trimer\_chain\_spiegato.ipynb}}

Mirror di \texttt{vqe\_trimer\_ring\_spiegato.ipynb}, con un adattamento fisico non
banale: il blocco a basso campo della catena ($B$) non è banale come il blocco $C$
usato nel punto di lavoro standard dell'anello, quindi la sezione di derivazione a
mano ($X,H,\mathrm{CX},X$) non si applica — sostituita con una verifica numerica di
\texttt{prepare\_state}.

\textbf{Verifiche e correzione.} Eseguito per intero: trovato un problema reale, non
un bug di codice ma di \textbf{presentazione} — due celle (VQE a un punto singolo, e
confronto su griglia) mostravano $\mathcal F=0.826612$ invece di $1.000000$, non per un
limite strutturale ma per un minimo locale specifico del seed di default (42).
Diagnosticato con più restart ($R=8$ risolve) e con seed diversi (qualunque seed
$\neq42$ risolve già a $R=4$). Corretto cambiando il seed esplicito nelle due celle;
rieseguito, ora $\mathcal F=1.000000$ esatto ovunque.

\textbf{Conclusione.} Un esempio concreto, verificato sul campo, della distinzione fra
"limite di espressività" e "artefatto di ottimizzazione" che il progetto tiene ferma
fin dal dimero: qui il rischio di scambiare l'uno per l'altro era reale e si è
materializzato, non solo teorico.

\section*{\texttt{vqe\_trimer\_chain\_W\_spiegato.ipynb}}

Mirror di \texttt{vqe\_trimer\_ring\_W\_spiegato.ipynb}, con lo stesso adattamento
della sezione di derivazione del notebook precedente (blocco $B$ non banale).

\textbf{Verifiche.} Eseguito per intero, zero errori: overlap di preparazione
$=1.000000$, HA e PMA entrambi $\mathcal F=1$ esatto sulla griglia ridotta.

\section*{Conclusione generale}

La catena ha ora una famiglia di file per la fase VQE senza DM parallela a quella
dell'anello (produzione RBS, produzione $W$, produzione senza branching, due guide
interattive), con lo stesso livello di verifica end-to-end. Due punti di attenzione
emersi durante la costruzione, entrambi reali e non solo teorici:

\begin{itemize}
\item \textbf{Il seed di default può ingannare}: il problema trovato in
\texttt{vqe\_trimer\_chain\_spiegato.ipynb} mostra che un valore di fidelity
"sospetto" (vicino a un tetto già noto per un'altra famiglia) va sempre verificato con
più restart o seed diversi prima di essere interpretato come limite strutturale.
\item \textbf{Ripetere un blocco non è come aggiungere parametri a un blocco fisso}:
la scoperta su \texttt{vqe\_trimer\_chain\_nobranch.py} (K=2 supera il tetto di K=1)
generalizza un principio già visto per l'anello sotto DM, qui confermato anche a
$D=0$.
\end{itemize}

\textbf{Nota.} \texttt{\seqsplit{confronto\_ansatz\_entangler\_trimero\_catena.ipynb}} non è
descritto in questo documento: pur essendo attualmente $D=0$-only, è lo stesso slot
di notebook che per l'anello è poi diventato il confronto sistematico esteso al DM
(\texttt{\seqsplit{confronto\_ansatz\_entangler\_trimero\_anello.ipynb}}). Verrà documentato,
nella sua forma definitiva, insieme al resto della Fase 5 (VQE con DM) per la catena,
non ancora affrontata.

\end{document}
